% !PRINTED 644
\weberpage{644}
Ist aber $\mathfrak{p}=\mathfrak{p}'$, so ist $\mathfrak{P}=\mathfrak{p}$, und die Normbildung ergiebt $2f_1=f$.

Welcher dieser beiden F\"alle eintritt, das h\"angt also davon ab, ob das Primideal $\mathfrak{p}$ die Substitution $(r,r^{-1})$ gestattet oder nicht, d.\,h. ob $(r,r^{-1})$ in der Gruppe, zu der das Primideal geh\"ort, enthalten ist oder nicht. Nach \S.\,172 ist dieser Unterschied dadurch bedingt, ob es irgend einen Exponenten $h$ giebt, f\"ur den die Congruenz
\[
p^h \equiv -1 \ (\mathrm{mod}\ m)
\]
erf\"ullt ist, oder ob es keinen solchen Exponenten giebt.

Wir fassen das Ergebniss folgendermaassen zusammen:

2. \gesperrt{Die von $q$ verschiedenen nat\"urlichen Primzahlen $p$ zerfallen in zwei Arten $p_1, p_2$.}

\gesperrt{Die erste Art $p_1$ ist dadurch charakterisirt, dass f\"ur irgend einen Exponenten $h$}
\[
p_1^h \equiv -1 \ (\mathrm{mod}\ m),
\]
\gesperrt{und diese Primzahlen zerfallen, wenn sie f\"ur den Modul $m$ zum Exponenten $f$ geh\"oren und $\mu=ef$ gesetzt wird, im K\"orper $H_m$ in $e$ Primfactoren $\frac12 f^{\text{ten}}$ Grades.}

\gesperrt{Ihre Primfactoren sind auch in $\Omega_m$ nicht weiter zerlegbar.}

\gesperrt{Die Primzahlen der zweiten Art $p_2$ sind dadurch bestimmt, dass keine ihrer Potenzen nach dem Modul $m$ mit $-1$ congruent wird, und sie zerfallen im K\"orper $H_m$ in $\frac12 e$ Primfactoren $f^{\text{ten}}$ Grades. Ihre Primfactoren sind in $\Omega_m$ noch in je zwei Factoren zerlegbar.}

Bei den Primzahlen der ersten Art muss $f$ sicher eine gerade Zahl sein. Wenn $m$ ungerade ist, so gen\"ugt es auch, damit $p$ eine Primzahl von der ersten Art sei, dass sie zu einem geraden Exponenten geh\"ore; denn dann ist
\[
(p^{\frac12 f}-1)(p^{\frac12 f}+1)\equiv 0\ (\mathrm{mod}\ m).
\]
Der erste dieser Factoren $p^{\frac12 f}-1$ ist aber nicht durch $m$ theilbar, weil sonst $p$ zum Exponenten $\frac12 f$ geh\"oren w\"urde. Folglich muss $p^{\frac12 f}+1$ durch $q$ theilbar sein. Hier k\"onnen aber nicht beide Factoren $p^{\frac12 f}-1$ und $p^{\frac12 f}+1$ durch $q$ theilbar sein, weil sonst auch ihre Differenz, die $=2$ ist, durch $q$ theilbar w\"are, und folglich ist $p^{\frac12 f}+1$ durch $m$ theilbar.
